\section{Research Thrusts}
\label{sec:thrusts}

The proposed work is organized into three research thrusts that
correspond to the major technical challenges underlying \sysname:
\emph{capturing} an open-source critical-infrastructure stack's
configuration and behavior
(Thrust~1, \S\ref{sec:thrust1}), \emph{replicating} it as a
controlled-fidelity digital twin (Thrust~2, \S\ref{sec:thrust2}), and
\emph{analyzing} that twin through agentic red teaming
(Thrust~3, \S\ref{sec:thrust3}). Although the thrusts are presented in
this order---and the timeline in \S\ref{sec:timeline} phases them
accordingly---they are deeply coupled: fidelity requirements from
Thrust~3 drive replication decisions in Thrust~2, which in turn drive
observation priorities in Thrust~1. Each thrust is accompanied by
specific research questions, the techniques we propose to investigate,
and success criteria.

%% =============================================================
\subsection{Thrust 1: Capturing Open-Source Stack Configurations}
\label{sec:thrust1}

\subsubsection*{Research questions}
\textbf{(RQ1.1)} What is the right representation for a living,
partially-known model of an open-source critical-infrastructure
deployment, such that AI agents can reason about it, update it
incrementally, and flag contradictions?
\textbf{(RQ1.2)} What observation techniques yield the highest
information per unit of operational risk when applied to a production
\cps?
\textbf{(RQ1.3)} How should static documentation and dynamic
observations be fused when they disagree, and when should the system
escalate to a human rather than resolving the disagreement itself?

\subsubsection{A Hybrid Knowledge Base}

A recurring observation in our preliminary work is that neither a
purely-structured representation nor a purely-unstructured one is
sufficient to capture an \cps.  A graph database cleanly expresses
``database node A is connected to application server B over subnet
S via service-mesh proxy X'' but awkwardly expresses ``the upstream
release notes observe, in the changelog for version 14.2, that
configurations with \texttt{synchronous\_commit=off} and streaming
replication may silently lose committed transactions under specific
failover conditions.''  A vector database cleanly expresses the
latter but poorly supports topological queries.  We therefore propose
a \emph{hybrid} knowledge base with two coordinated stores:

\begin{itemize}[leftmargin=1.3em,itemsep=0.1em,topsep=0.2em]
  \item A property-graph database (prototyped in Neo4j) storing
  \emph{entities} (hosts, services, devices, users, data stores,
  network segments, policies), \emph{relationships} (connects-to,
  depends-on, authenticates-via, trusts, monitors), and
  \emph{properties} (versions, CPEs, ownership, confidence, provenance).
  The schema is closed and curated; new node and edge types are
  introduced through a controlled process rather than by agents on
  their own authority.
  \item A vector database (prototyped in Qdrant or Weaviate) storing
  natural-language content---manual sections, configuration snippets,
  policy text, log samples, ticket narratives---each with a pointer
  back to the graph entity or relationship it describes, and a
  timestamp and source URI for provenance.
\end{itemize}

The two stores are linked by \emph{anchor edges}: every vector entry
has one or more anchor edges into the graph, and every graph entity
can enumerate the vector entries that reference it. Agent queries
typically traverse both: ``given this database cluster's engine and
version, what do the upstream release notes and deployment guides say
about authentication defaults, and does the network topology around
this cluster contradict those defaults?'' We will investigate whether a
structured prompting convention---retrieval over the vector store
constrained by a graph-traversal filter---yields better agent reasoning
than either store alone.

\paragraph{Research sub-problem: confidence and contradiction.}
Unlike a configuration-management database, which assumes its contents
are true, the \sysname KB must represent uncertainty and disagreement.
We propose to extend each fact with a \emph{confidence score}
(initially derived from source type: written documentation and project
READMEs are trusted less than observed network behavior, which is
trusted less than cryptographic attestation of running binaries), an
\emph{agreement set} (the observations that support the fact), and a
\emph{disagreement set} (the observations that contradict it). An open
research question is how to calibrate these scores in a way that
downstream agents---particularly Red Team Agents, whose attack plans
may hinge on a single load-bearing fact---can use.

\subsubsection{Observation Techniques}

Observer Agents gather information about the running \cps through
\emph{strictly passive} means. We propose to investigate five
complementary techniques, each with distinct coverage, cost, and
deployability profiles.

\paragraph{Passive network analysis.}
Observer Agents consume packet captures from network taps or SPAN
ports, using protocol dissectors to reconstruct flows and sessions.
For the application-layer protocols characteristic of open-source
stacks---database wire protocols (PostgreSQL, MySQL, MongoDB), message
broker protocols (Kafka, AMQP, NATS), RPC frameworks (gRPC,
Thrift), coordination protocols (etcd's Raft-over-gRPC,
ZooKeeper's ZAB), and cluster-internal HTTP---we will extend
existing open-source dissectors where available and build LLM-assisted
dissectors for custom or less-standardized variants, using the project
documentation stored in the KB as few-shot context. A key research
question is \emph{flow attribution}: given a flow between two hosts,
can the Observer infer the application-layer purpose (e.g.,
``application tier polling the authorization service once per request,
cache miss rate rising'') rather than just the protocol?

\paragraph{Log monitoring.}
Open-source infrastructure stacks produce logs in radically
inconsistent formats: structured JSON from modern services, line-oriented
syslog-like records from legacy Unix daemons, the idiosyncratic
text formats of widely-used servers (nginx access logs, PostgreSQL
logs, Kafka broker logs), and free-form application output. We
propose an LLM-based log normalization pipeline that
maps arbitrary log records into a shared ontology (event type,
source, target, outcome, attributes) rooted in the KB's graph schema.
The research contribution here is not the use of LLMs for
log parsing---that is now standard---but the closed-loop use of the KB
as grounding: because the agent knows what hosts and services exist,
it can validate parses against plausible entities and flag spurious
matches.

\paragraph{API monitoring.}
Modern \cps deployments increasingly use microservice-style
architectures throughout the stack. Where this is the case, we
will use service-mesh facilities (e.g., Istio with Envoy or Linkerd)
or eBPF-based probes to capture inter-service API calls without
touching application code. This gives the Observer access to
semantics---which service called which API on which resource---that
network-layer observation cannot.

\paragraph{Configuration monitoring.}
Open-source critical-infrastructure stacks are typically managed
through configuration managers such as Ansible, Puppet, Chef, or
SaltStack, through infrastructure-as-code systems such as Terraform
and Helm, and through GitOps workflows (e.g., ArgoCD, Flux) that
continuously reconcile Kubernetes clusters with versioned manifests.
Where such a system is present, an Observer Agent with read-only
access to its inventory, manifests, and playbooks can extract a large
fraction of the system's intended configuration directly, without ever
touching a production host. We will build adapters for the major
open-source configuration and GitOps systems and, where third-party
vendors provide proprietary management APIs, evaluate the cost of
extending the approach.

\paragraph{Security-tool integration.}
Finally, Observer Agents will integrate with existing security tooling
when it is present: endpoint detection and response (EDR) agents,
cloud-native runtime security products (e.g., those in the
Falco/Wiz/Sysdig category), container-image scanners, and security
information and event management (SIEM) systems. The value here is
not to replace these tools but to treat their alerts as
observations---signals that the KB should record and that Red Team
Agents should consider when modeling what an attacker would see.

\paragraph{Research sub-problem: observation planning as an optimization
problem.}
Observer Agents do not enjoy unlimited visibility. SPAN ports are
scarce; storage and bandwidth are finite; some network segments can
only be observed through a single, contended vantage point. We propose
to formulate observation planning as a constrained optimization: given
a set of KB entities of interest, a set of possible vantage points,
and a budget, maximize the expected information gain about the
uncertain parts of the KB. We will investigate whether this problem
admits a tractable approximation and whether solutions computed by a
classical optimizer can be improved by LLM-guided reasoning about
which facts are security-relevant.

\subsubsection{Expected Outcomes of Thrust~1}
\begin{itemize}[leftmargin=1.3em,itemsep=0.1em,topsep=0.2em]
  \item A schema and open-source implementation of a hybrid,
  uncertainty-aware knowledge base for open-source
  critical-infrastructure stacks.
  \item An extensible library of passive observation modules, including
  adapters for common open-source application-layer protocols, log
  normalizers, service-mesh probes, configuration-manager and GitOps
  readers, and SIEM connectors.
  \item An evaluation, on reference \CPSs (see
  \S\ref{sec:evaluation}), of the coverage, latency, and
  operational-risk profile of the above techniques.
\end{itemize}

%% =============================================================
\subsection{Thrust 2: Replicating Deployments}
\label{sec:thrust2}

\subsubsection*{Research questions}
\textbf{(RQ2.1)} What does it mean for a digital twin to be
``sufficiently faithful'' for security analysis, and how can that
property be measured?
\textbf{(RQ2.2)} Given that a perfectly faithful twin of a large \cps
is infeasible, how should a replication algorithm trade fidelity
against cost, and what parameters should the operator be able to
control?
\textbf{(RQ2.3)} For components that cannot be reproduced directly
in the twin, how can an observed behavior be turned into a mock that
is sound for security analysis?

\subsubsection{Controlled-Fidelity Replication}

The naive approach to digital twin construction---replicate every
component---fails for three reasons: large \CPSs contain dozens or
hundreds of service instances, many of which are replicas of the
same underlying binary with only small configuration differences;
some components in the stack are closed-source or depend on
third-party services that cannot be copied wholesale into the twin;
and the operational cost of running a full-scale replica (especially
of data tiers whose state is measured in terabytes) often exceeds the
cost of the original. \sysname's replication approach is therefore
explicitly \emph{representational} rather than \emph{exhaustive}: the
twin preserves the features that matter for security analysis and
deliberately abstracts away features that do not.

We formalize this as a \emph{controlled-fidelity replication problem}.
Given the KB's inventory of components, the twin is characterized by
three parameters:

\begin{description}[leftmargin=1.6em,style=nextline,itemsep=0.15em,topsep=0.3em]
  \item[Diversity ($D$).] Every equivalence class of components in
  the original system (defined by project, version, configuration
  profile, and role---e.g., ``PostgreSQL 14.8 primary with
  streaming replication enabled and \texttt{pg\_hba} configured for
  LDAP auth'') must be represented by at least one instance in the
  twin. This is the minimum condition for the Red Team to encounter
  every version- and configuration-specific bug class present in the
  real system.
  \item[Topology representation ($T$).] Many vulnerabilities depend
  on \emph{patterns} of connectivity rather than individual links
  (e.g., ``application tier and database tier share a flat network
  segment with no egress filtering'', ``the secrets manager is
  reachable from the build cluster via a single overly-permissive
  IAM policy''). Our algorithm enumerates recurring
  topological patterns in the original system and guarantees that
  each pattern appears at least once in the twin.
  \item[Data representation ($R$).] Some attacks depend on the data
  being processed (e.g., SQL injection attacks that depend on
  particular column types and constraints; deserialization attacks
  that depend on specific class paths being available; business-logic
  attacks that depend on realistic account, role, or tenant
  distributions). The twin is populated with synthetic
  data whose statistical and schema-level properties match the
  original, under a privacy constraint that excludes
  personally-identifiable and safety-sensitive records.
\end{description}

Given $(D, T, R)$, the replication algorithm solves a set-cover-like
problem: select a minimum-size set of component instances that
covers all equivalence classes, all required topological patterns,
and all required data-distribution targets. We will investigate
approximation algorithms for this problem and develop an
operator-facing interface that exposes $(D, T, R)$ as tunable knobs,
with default settings derived from the KB's risk model.

\paragraph{Research sub-problem: quantifying fidelity.}
A central open question is how to measure whether a twin is
``faithful enough.'' We propose two complementary metrics. The first
is \emph{attack reproducibility}: for a seeded set of known-real
vulnerabilities (injected deliberately into the real system and
documented in the KB, or drawn from historical CVE reports on identical
component versions), what fraction can the Red Team reproduce against
the twin? The second is \emph{behavioral divergence}: for a stream of
observations on the real system, how often does the twin produce a
statistically distinguishable response to the same stimulus? We will
investigate both metrics empirically on reference testbeds.

\subsubsection{Replicating Opaque and Hard-to-Virtualize Components}

A substantial fraction of the components in a typical
critical-infrastructure deployment cannot be reproduced by simply
spinning up another container from an upstream image. These include
closed-source third-party services that the open-source stack depends
on (identity providers, payment gateways, vendor-supplied connectors),
services whose state is too large or too sensitive to copy into the
twin (production databases, object stores), and custom components
whose source is not available to the assessment team. Our group has
spent the last decade developing techniques to model exactly these
kinds of opaque systems---techniques that transfer naturally to
software-component replication, as described below. \sysname builds
directly on that foundation.

\paragraph{Direct replication from upstream artifacts.}
For the great majority of components in an open-source stack, the
Replication Agent constructs a twin instance by pulling the same
upstream artifact (container image, binary package, or source build)
that production uses, pinned to the same version, and applying the
configuration recovered from the KB. This is the common case for
open-source software and is what makes \sysname's approach tractable
in ways that would be impossible for proprietary stacks.

\paragraph{Configuration and state sanitization.}
For stateful components---databases, object stores, message brokers
with durable queues---the twin is built from the same binary as
production but populated with a sanitized replica of state: schemas
and indexes are preserved, row counts are preserved by percentile, but
personally-identifiable and safety-sensitive records are replaced with
synthetic equivalents that match the distribution of the originals.
Our prior work on HALucinator [Clements et al., USENIX Security 2020],
which automatically substitutes high-level implementations for
low-level components whose faithful reproduction is infeasible, provides
a methodological template: rather than replacing hardware-abstraction-layer
calls, the Replication Agent here replaces specific rows, files, or
message payloads with synthetic analogues while preserving the
surrounding structure.

\paragraph{Behavioral mock synthesis.}
For components where direct replication is infeasible---closed-source
third-party services, external APIs, or components whose dependencies
cannot be satisfied in the twin---we will apply ideas from our
Conware [Spensky et al., AsiaCCS 2021] work, which automatically models
opaque systems from recorded interactions. A Replication Agent will
construct a finite-state model of the component's externally visible
behavior from Observer traces, and instantiate that model as a mock
service that responds to queries as the real component would. Mocks
are not a substitute for real execution, and we are careful to label
findings produced against them as such; but they allow the surrounding
system to be reasoned about without the mocked component becoming a
hole that disables all analysis.

\paragraph{Human escalation for irreducible gaps.}
Some components will resist all of the above approaches: the real
component depends on credentials or data that the assessment team is
not authorized to replicate; the service is itself an operator
(e.g., a human approval step in a workflow); or the component
performs an external side-effect (sending email, billing a customer)
whose simulation raises its own concerns. In these cases, the
Replication Agent surfaces a structured request to the human operator,
who may provide additional information, authorize an offline
reconstruction, or declare the component out of scope for the current
twin.

\subsubsection{Expected Outcomes of Thrust~2}
\begin{itemize}[leftmargin=1.3em,itemsep=0.1em,topsep=0.2em]
  \item A formalization of the controlled-fidelity replication problem
  and one or more approximation algorithms.
  \item Metrics for quantifying twin fidelity, together with an
  empirical evaluation on reference \CPSs.
  \item Open-source tooling that drives agent-directed twin
  construction from upstream artifacts, applying ideas adapted from
  HALucinator and Conware to the replacement and behavioral modeling
  of opaque software components.
\end{itemize}

%% =============================================================
\subsection{Thrust 3: Analyzing Deployments}
\label{sec:thrust3}

\subsubsection*{Research questions}
\textbf{(RQ3.1)} How should component-level vulnerability analyses
(fuzzing, static analysis, symbolic execution) be scheduled and
composed against a digital twin of non-trivial size?
\textbf{(RQ3.2)} Can multi-step attack plans generated by AI-planning
techniques and refined by LLM-based reasoning discover vulnerabilities
that component-level analyses miss?
\textbf{(RQ3.3)} How should an agentic red team triage its own
findings before escalating to a human, so that the human's time is
spent on attacks that are both likely real and likely
high-consequence?

\subsubsection{Component-Level Analyses}

The twin is, by design, a controllable and repeatable target. This
allows a wide range of analyses that would be unsafe against the real
system:

\begin{itemize}[leftmargin=1.3em,itemsep=0.1em,topsep=0.2em]
  \item \emph{Coverage-guided fuzzing} of services, using tools such
  as AFL++\ and LibAFL against network-accessible components, including
  application-protocol fuzzers seeded with corpora drawn from Observer
  traffic captures.
  \item \emph{Static analysis} of service binaries and, where source
  is available, source code, using the \texttt{angr} binary-analysis
  framework (co-developed by the PIs) and LLM-assisted bug-pattern
  detectors that specialize in authentication handlers, deserializers,
  and request-routing logic.
  \item \emph{Symbolic execution} for authentication and authorization
  logic, where the state space is small enough to explore exhaustively.
  \item \emph{Configuration analysis} that evaluates the twin's
  configuration against a library of known misconfiguration patterns
  (default credentials, permissive network policies, insecure
  cryptographic settings, and deployment-specific patterns such as
  ``build-system service account with cluster-admin on the production
  namespace'' or ``secrets manager reachable from a publicly-exposed
  ingress controller'').
\end{itemize}

Each component-level analysis produces \emph{atomic findings}: concrete
weaknesses characterized by a precondition (what an attacker must
already have in order to exploit the weakness) and a postcondition
(what the attacker gains). Atomic findings are pushed into the KB,
where they become available to the system-level planner.

\subsubsection{System-Level Multi-Step Attack Planning}

The central intellectual claim of Thrust~3 is that individual atomic
findings, while useful, are not where the most severe \cps
vulnerabilities live. Real incidents---the SolarWinds supply-chain
compromise, the Log4Shell exploitation wave, the MOVEit data-theft
campaign, the Colonial Pipeline ransomware event---invariably involve
\emph{chains}: initial access through one weakness, lateral movement
through another, privilege escalation through a third, and
consequential action through a fourth. Any automated red team that
stops at atomic findings will systematically miss the class of attacks
operators most need to know about.

Our ChainReactor work [Pasquale et al., USENIX Security 2024] demonstrated that AI
planning---specifically, PDDL-based planning over atomic attacks
characterized by preconditions and postconditions---can discover
privilege-escalation chains in Unix environments that even expert
red-teamers missed. \sysname extends this approach in three directions.

\paragraph{Domain expansion.}
Where ChainReactor operated on a single host, \sysname operates on a
multi-host, multi-service, deployment-wide domain. We will extend the
PDDL formulation to model network segments, service-to-service trust
relationships, identity- and token-scoped attacker capabilities, and
data-tier effects (e.g., ``attacker can read rows from table $T$'',
``attacker can write to topic $K$'', ``attacker holds a valid token
for service $S$''). The state-space explosion that results will be
managed by a combination of hierarchical planning (plan at the
segment or namespace level, then refine within) and LLM-guided
heuristic search, in which a language model proposes candidate next
steps that a classical planner then validates.

\paragraph{Grounded plan synthesis with LLMs.}
Recent evaluations show that AI agents are making measurable progress
on multi-step attack scenarios, and that even relatively small models
can propose plausible attack steps when given appropriate context.
However, LLMs are also known to hallucinate preconditions and
postconditions, which is unacceptable for an automated red team. We
will investigate a hybrid approach in which LLMs propose candidate
attack plans in natural language, a grounding module converts each
step into PDDL actions (refusing steps that do not correspond to
atomic findings in the KB), and a classical planner validates the
result. Our recent experience with the Artiphishell system in the
DARPA AIxCC competition provides a working foundation for this style
of LLM/classical-planner hybrid.

\paragraph{Attack execution and confirmation.}
A plan that is valid on paper is not yet evidence of a vulnerability.
\sysname instantiates the twin, executes the plan against it using a
library of attack primitives implemented as Python agents, and checks
whether the nominated flag is reached. Successful plans are replayed
on a clean twin snapshot to eliminate side-effects introduced by
earlier attempts. Failed plans are examined for their failure reasons,
which are written back into the KB and used to refine future planning.

\subsubsection{Findings Triage and Human Escalation}

An unfiltered agentic red team will generate more findings than a
human team can validate. \sysname therefore includes a triage stage
that ranks findings before presenting them. The ranking draws on four
signals: the \emph{confidence} with which the KB entities supporting
the finding are grounded in observations (as opposed to documentation
alone); the \emph{consequence class} of the compromised flag (as
defined in the operator's security policy); the \emph{path length} of
the attack (shorter chains are more likely to be realistic); and the
\emph{novelty} of the finding relative to previously-reported ones.
Triage is itself an agentic process, and its outputs are auditable:
the analyst receives a ranked list with a short natural-language
justification for each item, and can drill into the provenance trace
that supports it.

\subsubsection{False Positives and the Fidelity Feedback Loop}

When a finding is labeled as a false positive by the human analyst,
the cause is almost always one of two things: a divergence between
the twin and the real system (``that port isn't actually open in
production''), or a policy misalignment (``that user does have
legitimate access to that resource''). Both cases produce structured
feedback into the KB. Divergences trigger the Observer Agents to
re-investigate the relevant corner of the system; policy
misalignments trigger a KB update that Red Team Agents will consider
in future planning. This closed loop is what distinguishes \sysname
from a one-shot penetration test: over time, both the twin's fidelity
and the red team's precision should improve.

\subsubsection{Expected Outcomes of Thrust~3}
\begin{itemize}[leftmargin=1.3em,itemsep=0.1em,topsep=0.2em]
  \item A PDDL-based attack-planning domain for multi-host, multi-service
  open-source deployments, together with a hybrid LLM/classical planner.
  \item An open-source agentic red-team framework that composes
  component-level analyses into end-to-end attack plans, executable
  against \sysname digital twins.
  \item An empirical study on reference \CPSs of the precision, recall,
  and path-length distribution of discovered attacks, with comparison
  to human-red-team baselines where available.
\end{itemize}
